\section{Impact}
\label{sec:Impact}
The IoT-Sim provides an open and modular experimental environment for studying and evaluating attack detection models in IoT networks. Its client-side implementation makes it easily accessible and extensible, lowering the barrier to reproducing IoT scenarios and testing AI-based intrusion detection systems without the need for specialized hardware or complex deployments.

The simulator was used in a controlled laboratory environment to evaluate two attack detection models trained under different paradigms, centralized and federated learning~\cite{federated_learning}, demonstrating its suitability for comparing heterogeneous training strategies under identical network conditions. By enabling reproducible and configurable experiments, it facilitates the systematic analysis of model robustness, adaptability, and performance in diverse attack scenarios.

Beyond these initial experiments, the IoT-Sim opens up new research directions related to lightweight and distributed detection mechanisms for resource-constrained IoT devices. It also allows researchers to analyze model behavior under different network configurations and attack conditions, thereby supporting the evaluation of detection performance in diverse IoT environments.

Although its current use remains within the research laboratory, the simulator holds strong potential for broader adoption in academic and industrial contexts. Future developments will focus on expanding the network simulation capabilities and incorporating new communication protocols, traffic models, and attack behaviors to achieve more realistic and versatile experimentation environments for intelligent IoT network management.
